\chapter{GitOps Pipeline and CI/CD}
\label{ch:workflow}

Chapter~\ref{ch:kubernetes} described the Kubernetes infrastructure---Crossplane XRDs, Argo Events, Argo Workflows---that implements the quality gate system. This chapter describes how that infrastructure operates within the broader GitOps pipeline: the promotion-driven quality gate flow, the event-driven validation pipeline end-to-end, the single-repo architecture and its ``Bring Your Own CI'' design principle, the GitOps Promoter's multi-environment model, human-in-the-loop decision patterns, and the quality gate as a closed-loop control system. The chapter also describes the three CI/CD pipelines that build, test, and release the DriveBy system itself.

Together, Chapters~\ref{ch:architecture},~\ref{ch:kubernetes}, and this chapter form a three-layer value chain: the DriveBy CLI provides the validation logic, Kubernetes infrastructure orchestrates its execution, and GitOps pipelines make validation results drive promotion decisions. This chapter describes the outermost layer---the point where validation results become operational signals that advance or block software delivery.

The chapter addresses Research Question~3 (RQ3): \textit{How effectively does specification-driven quality assurance integrate into cloud-native GitOps pipelines as an autonomous quality gate?} The answer is demonstrated through a production deployment on the \texttt{private.novelcore.org} cluster, where DriveBy validates a reference API across three environments.

% ============================================================================
\section{Promotion-Driven Quality Gate}
\label{sec:promotion-driven-gate}
% ============================================================================

The quality gate operates on a key insight: \textit{validate the source environment, not the target}. When a developer merges code into the application repository, the CI pipeline builds and deploys the new version to the development environment. ArgoCD~\cite{argocd} synchronizes the development environment automatically. The GitOps Promoter detects the successful synchronization and creates a promotion pull request proposing to update the staging environment's manifest to the new version.

This promotion PR triggers the quality gate. The GitHub webhook fires, the EventSource receives the event, the Sensor filters for PR actions, and the Workflow executes validation against the \textit{development} environment---the environment where the new code is already running. The validation results determine whether the promotion PR is approved (merged) or blocked (left open with a failure status).

This inversion is deliberate. Validating the source environment ensures that quality is assessed against the actual deployment rather than a hypothetical future state. If the API in development passes strict DDT validation, the same code deployed to staging will behave identically (assuming equivalent infrastructure). The alternative---validating the target environment after promotion---would detect problems only after the potentially breaking change has already been deployed, violating the shift-left principle~\cite{shift-left}.

Figure~\ref{fig:gitops-pipeline} shows the end-to-end pipeline from developer pull request to validation feedback.

\begin{figure}[htbp]
\centering
\resizebox{\textwidth}{!}{%
% GitOps Event-Driven Pipeline with Source Hydrator — Chapter 6
% Developer commits dry manifests → Hydrator renders → -next branch → Promoter PR → Webhook → Argo Events → Workflow → Feedback
\begin{tikzpicture}[
    >={Stealth[length=2.5mm]},
    node distance=0.6cm,
    stage/.style={
        rectangle,
        rounded corners=3pt,
        draw=#1!60!black,
        fill=#1!12,
        minimum width=2cm,
        minimum height=0.8cm,
        font=\scriptsize,
        align=center
    },
    lane label/.style={
        font=\scriptsize\bfseries,
        rotate=90,
        anchor=center,
        text=#1!70!black
    },
    lane bg/.style={
        fill=#1!5,
        draw=#1!20,
        rounded corners=4pt
    }
]

% Lane backgrounds
\begin{scope}[on background layer]
    \node[lane bg=red,    minimum width=2.4cm, minimum height=3.5cm] at (-2,0) {};
    \node[lane bg=violet, minimum width=2.4cm, minimum height=3.5cm] at (0.3,0) {};
    \node[lane bg=blue,   minimum width=3.2cm, minimum height=3.5cm] at (2.8,0) {};
    \node[lane bg=cyan,   minimum width=2cm, minimum height=3.5cm] at (5,0) {};
    \node[lane bg=teal,   minimum width=3.6cm, minimum height=3.5cm] at (7.4,0) {};
    \node[lane bg=orange, minimum width=4cm, minimum height=3.5cm] at (11,0) {};
\end{scope}

% Lane labels (inside the lane background, near the top edge of each rectangle)
% Round-10 follow-up to supervisor [408]: previously sat outside the lane bg at y=1.9
% (above the rectangle); now positioned at y=1.4, inside each lane background
% (which extends from y=-1.75 to y=+1.75) above the top row of stages at y=0.8.
\node[lane label=red]     at (-3.1,1.4)  {BYOCI};
\node[lane label=violet]  at (-0.8,1.4)  {Hydrator};
\node[lane label=blue]    at (1.3,1.4)   {GitHub};
\node[lane label=cyan]    at (4.1,1.4)   {Ingress};
\node[lane label=teal]    at (5.7,1.4)   {Argo Events};
\node[lane label=orange]  at (9.1,1.4)   {Argo Workflows};

% BYOCI lane (developer trigger)
\node[stage=red] (deploy) at (-2,0.8)  {Developer\\commits dry};
\node[stage=red] (main)   at (-2,-0.5) {\texttt{main}\\branch};
\draw[->,thick] (deploy) -- (main);

% Hydrator lane
\node[stage=violet] (hydrator) at (0.3,0.8)  {ArgoCD\\Hydrator};
\node[stage=violet] (nextbr)   at (0.3,-0.5) {\texttt{-next}\\branch};
\draw[->,thick] (main) -- (hydrator);
\draw[->,thick] (hydrator) -- (nextbr);

% GitHub lane
\node[stage=blue] (promoter) at (2.3,0.8)  {Promoter\\creates PR};
\node[stage=blue] (webhook)  at (2.3,-0.5) {GitHub\\Webhook};
\draw[->,thick] (nextbr) -- (promoter);
\draw[->,thick] (promoter) -- (webhook);

% Ingress lane
\node[stage=cyan] (traefik) at (5,0.1) {Traefik\\Ingress};
\draw[->,thick] (webhook) -- (traefik);

% Argo Events lane
\node[stage=teal] (es)     at (6.3,0.8)  {EventSource\\(webhook)};
\node[stage=teal] (eb)     at (6.3,-0.5) {EventBus\\(JetStream)};
\node[stage=teal] (sensor) at (8.5,0.1)  {Sensor\\(filter)};
\draw[->,thick] (traefik) -- (es);
\draw[->,thick] (es) -- (eb);
\draw[->,thick] (eb) -- (sensor);

% Argo Workflows lane
\node[stage=orange] (wf)      at (10.2,0.8)  {Argo\\Workflow};
\node[stage=orange] (status)  at (10.2,-0.5) {Commit\\Status};
\node[stage=orange] (comment) at (12,-0.5)   {PR\\Comment};
\draw[->,thick] (sensor) -- (wf);
\draw[->,thick] (wf) -- (status);
\draw[->,thick] (wf) -- (comment);

% Feedback loop back to Promoter PR
\draw[->,thick,blue!60!black,dashed] (status.south) -- ++(0,-0.6) -| ($(promoter.south)+(0.3,-0.6)$) -- ($(promoter.south)+(0.3,0)$);

% BYOCI boundary annotation
\draw[dashed,thick,red!60!black] (-0.8,1.8) -- (-0.8,-1.8);
\node[font=\tiny\bfseries,text=red!60!black,anchor=north,rotate=90] at (-0.75,-0.5) {BYOCI boundary};

\end{tikzpicture}
%
}
\caption{End-to-end GitOps validation pipeline. A developer opens a promotion PR; the webhook triggers Argo Events, which submits an Argo Workflow. Validation results are posted back as commit status and PR comment (dashed feedback loop).}
\label{fig:gitops-pipeline}
\end{figure}

% ============================================================================
\section{Event-Driven Validation Pipeline}
\label{sec:event-pipeline-detail}
% ============================================================================

The event-driven pipeline transforms an external webhook into a structured validation workflow through five stages: ingress, event capture, event filtering, workflow submission, and feedback reporting. This section traces a single validation trigger through the entire pipeline.

% ----------------------------------------------------------------------------
\subsection{Webhook Delivery and Ingress}
\label{subsec:webhook-delivery}
% ----------------------------------------------------------------------------

The pipeline begins when GitHub delivers a webhook to the Ingress endpoint. The webhook URL follows the pattern:

\begin{verbatim}
https://{app}-{gate}-webhook.{domain}/{trigger-endpoint}
\end{verbatim}

For the reference deployment, this resolves to \texttt{https://perfect-api-staging-gate-webhook.private.novelcore.org/perfect-api-staging-gate-pr-trigger}. The Traefik Ingress controller terminates TLS (using a certificate issued by cert-manager), routes the request to the EventSource Service, and delivers the payload to the webhook listener on port~12000.

% ----------------------------------------------------------------------------
\subsection{Event Processing}
\label{subsec:event-processing}
% ----------------------------------------------------------------------------

The EventSource receives the webhook payload and publishes it to the EventBus (NATS JetStream). The Sensor subscribes to the EventBus and evaluates the event against its configured filters. For promotion PR validation, the Sensor checks that the event type is \texttt{pull\_request} and the action is one of \texttt{opened}, \texttt{reopened}, or \texttt{synchronize}. Events that do not match---such as PR label changes, review requests, or comment events---are silently discarded.

For matching events, the Sensor extracts ten parameters from the webhook payload using JSONPath expressions (detailed in Table~\ref{tab:sensor-params} of Chapter~\ref{ch:kubernetes}) and constructs a workflow submission request. The submission references the WorkflowTemplate generated by the \texttt{XSDLC} composition and passes the extracted parameters as workflow arguments.

% ----------------------------------------------------------------------------
\subsection{Latency Analysis}
\label{subsec:latency}
% ----------------------------------------------------------------------------

The end-to-end latency from webhook delivery to workflow completion depends on three factors: event processing overhead (typically under 2~seconds for EventSource $\rightarrow$ EventBus $\rightarrow$ Sensor $\rightarrow$ workflow submission), source readiness wait (0--120~seconds, depending on whether the deployment is already healthy), and validation execution time (10--60~seconds, depending on specification size and validation mode). In practice, a typical validation cycle for the reference API completes in under 90~seconds from PR creation to commit status posted.

% ============================================================================
\section{Single-Repo Architecture}
\label{sec:single-repo}
% ============================================================================

The quality gate system described in this thesis operates within a \textit{single-repository} model that deliberately separates application CI from platform-managed delivery. This section describes the architectural rationale, the boundary between developer and platform responsibilities, and the manifest propagation mechanism that bridges the two.

% ----------------------------------------------------------------------------
\subsection{Why a Single Repository}
\label{subsec:why-single-repo}
% ----------------------------------------------------------------------------

Traditional GitOps deployments maintain a separate \textit{config repository} alongside the application's source repository. The config repo holds Kubernetes manifests for each environment; a CI pipeline in the source repo updates the config repo on each build. This pattern doubles the number of repositories, fragments the change history across two Git logs, and introduces a synchronization pipeline (``update the config repo'') that must itself be maintained and monitored.

The XSDLC model eliminates this duplication. The application repository holds both source code and Kubernetes manifests in a single tree:

\begin{verbatim}
perfect-api/
  src/                     # Application source
  Dockerfile               # Build specification
  manifests/
    deployment.yaml        # Kubernetes deployment manifest
    service.yaml           # Kubernetes service manifest
\end{verbatim}

The developer maintains the manifests alongside the source code. When a CI pipeline builds a new container image, it updates \texttt{manifests/deployment.yaml} with the new image tag and pushes to \texttt{main}. The XSDLC delivery pipeline takes over from there, propagating manifests to environment branches, gating promotions, and syncing through ArgoCD. The developer never creates or manages environment branches, promotion PRs, or ArgoCD Applications---the XSDLC composition provisions all of these.

% ----------------------------------------------------------------------------
\subsection{Bring Your Own CI (BYOCI)}
\label{subsec:byoci}
% ----------------------------------------------------------------------------

The single-repo model embodies a deliberate design principle: \textbf{XSDLC owns delivery, not CI}. The boundary is explicit: \textit{XSDLC's responsibility begins when manifests land on \texttt{main}}.

Everything \textit{before} that boundary is the developer's responsibility: building the application, running unit tests, generating the container image, and updating the manifest with the new image tag. These tasks are inherently application-specific---a Python FastAPI service uses different build tools, test frameworks, and container patterns than a Go gRPC service. No platform abstraction can usefully generalize this diversity without either constraining developer choice or reimplementing CI functionality that mature tools (GitHub Actions, Jenkins, GitLab CI) already provide.

Everything \textit{after} that boundary is XSDLC's responsibility: propagating manifests to environment branches, triggering quality gates, managing promotion PRs, posting commit statuses, and syncing ArgoCD Applications. These tasks are platform-generic---the same propagation, gating, and promotion pattern applies regardless of the application's language, framework, or build tool.

This separation contrasts with opinionated CD platforms (Tekton Pipelines, Jenkins X) that attempt to own the full CI/CD lifecycle, providing build templates, test runners, and deployment strategies within a single framework. Such platforms trade flexibility for integration: they work well when the application fits their model and poorly when it does not. XSDLC is composable rather than monolithic---it integrates with any CI system that can push manifest changes to \texttt{main}.

% ----------------------------------------------------------------------------
\subsection{Manual Deploy Flow}
\label{subsec:manual-deploy}
% ----------------------------------------------------------------------------

The bridge between the developer and the XSDLC delivery pipeline is the \texttt{driveby-deploy.yml} workflow, generated by the XSDLC composition and committed to the application repository via \texttt{provider-upjet-github} (Section~\ref{subsec:generated-github}).

The workflow is triggered manually via \texttt{workflow\_dispatch} with three inputs: the source branch to take manifests from, the target environment (presented as a choice dropdown generated from the XSDLC's environments array), and the container image tag to deploy. When triggered, it:

\begin{enumerate}
    \item Checks out the \texttt{environment/<name>-next} branch (creating an orphan branch on first deploy).
    \item Cleans the branch root (removing previous manifest state).
    \item Copies the \texttt{manifests/} directory from the source branch to the branch root, flattening the directory structure.
    \item Stamps the image tag into \texttt{deployment.yaml} via \texttt{sed}.
    \item Commits and pushes to the \texttt{-next} branch.
\end{enumerate}

This design makes deployment explicit and targeted: the developer chooses which branch to deploy and which environment to target. Unlike an automated propagation model that fans out to all environments simultaneously, the manual trigger allows deploying a feature branch to \texttt{dev} without affecting \texttt{staging} or \texttt{prod}.

% ----------------------------------------------------------------------------
\subsection{Branch Topology}
\label{subsec:branch-topology}
% ----------------------------------------------------------------------------

Figure~\ref{fig:single-repo-branch-flow} shows the complete branch topology for a three-environment pipeline.

\begin{figure}[htbp]
\centering
\resizebox{\textwidth}{!}{%
% Two-Repo Branch Flow with Source Hydrator — Chapter 6, Section 6.3.3
% Developer commits dry manifests on main → hydrator renders per-env overlays → env-next → Promoter PRs → env → ArgoCD
\begin{tikzpicture}[
    >={Stealth[length=2.5mm]},
    branch/.style={
        rectangle,
        rounded corners=3pt,
        draw=#1!60!black,
        fill=#1!12,
        minimum width=3.2cm,
        minimum height=0.7cm,
        font=\scriptsize,
        align=center
    },
    action/.style={
        rectangle,
        rounded corners=2pt,
        draw=gray!50,
        fill=gray!5,
        minimum width=2.8cm,
        minimum height=0.5cm,
        font=\tiny,
        align=center
    },
    boundary/.style={
        dashed,
        thick,
        #1!60!black
    }
]

% Developer commit on main (top)
\node[action, minimum width=5cm] (push) at (5,0) {Developer commits dry manifests\\on \texttt{main} branch};

% BYOCI boundary line and side-anchored label
% The dashed boundary separates the developer's side (above) from the
% XSDLC platform side (below). One single label is anchored at the
% right edge, positioned to the right of the box stack so it does not
% collide with any other element.
\draw[boundary=red!70!black] (-0.5,-0.85) -- (10.5,-0.85);
\node[font=\tiny\bfseries,text=red!60!black,anchor=west,align=left]
  at (10.6,-0.85)
  {BYOCI\\boundary};

% ArgoCD Hydrator step
\node[branch=violet, minimum width=8cm] (hydrator) at (5,-1.7) {ArgoCD Source Hydrator renders per-environment overlays};

\draw[->,thick,violet!50!black] (push.south) -- (hydrator.north);

% env-next branches in gitops repo
\node[branch=blue] (devnext)     at (0,-3.0) {environment/\\dev-next};
\node[branch=blue] (stagingnext) at (5,-3.0) {environment/\\staging-next};
\node[branch=blue] (prodnext)    at (10,-3.0) {environment/\\prod-next};

\draw[->,thick,blue!50!black] (hydrator.south) -- ++(0,-0.1) -| (devnext.north);
\draw[->,thick,blue!50!black] (hydrator.south) -- (stagingnext.north);
\draw[->,thick,blue!50!black] (hydrator.south) -- ++(0,-0.1) -| (prodnext.north);

% Label showing hydrated output — anchored next to the centred
% staging-next arrow so it sits over empty space rather than crossing
% other arrows or boxes.
\node[font=\tiny,text=blue!60!black,anchor=west] at (5.3,-2.4) {manifests/ + hydrator.metadata};

% Promoter PRs
\node[action] (pr1) at (0,-4.0) {auto-merge (no gate)};
\node[action] (pr2) at (5,-4.0) {Promoter PR};
\node[action] (pr3) at (10,-4.0) {Promoter PR};

\draw[->,thick,blue!40!black] (devnext) -- (pr1);
\draw[->,thick,blue!40!black] (stagingnext) -- (pr2);
\draw[->,thick,blue!40!black] (prodnext) -- (pr3);

% Quality gates
\node[branch=orange, minimum width=2cm] (gate1) at (5,-4.9) {SDT Gate};
\node[branch=orange, minimum width=2cm] (gate2) at (10,-4.9) {Load Gate};

\draw[->,thick,orange!50!black] (pr2) -- (gate1);
\draw[->,thick,orange!50!black] (pr3) -- (gate2);

% Active branches
% Round-2 review [264]: row moved from y=-5.8 to y=-6.2. The two-line branch
% nodes previously left <2.5mm between the gate's south border and the branch
% node's north border, so the gate1->staging arrow (head longer than the gap)
% rendered clipped. All arrows are now border-anchored with room to render
% completely, which also restores a correct bounding box.
\node[branch=teal] (dev)     at (0,-6.2) {environment/\\dev};
\node[branch=teal] (staging) at (5,-6.2) {environment/\\staging};
\node[branch=teal] (prod)    at (10,-6.2) {environment/\\prod};

\draw[->,thick,teal!50!black] (pr1.south) -- (dev.north);
\draw[->,thick,teal!50!black] (gate1.south) -- (staging.north);
\draw[->,thick,teal!50!black] (gate2.south) -- (prod.north);

% ArgoCD sync (all auto)
\node[branch=green!60!black] (argodev) at (0,-7.6) {ArgoCD sync\\(auto)};
\node[branch=green!60!black] (argostg) at (5,-7.6) {ArgoCD sync\\(auto)};
\node[branch=green!60!black] (argoprod) at (10,-7.6) {ArgoCD sync\\(auto)};

\draw[->,thick,green!40!black] (dev.south) -- (argodev.north);
\draw[->,thick,green!40!black] (staging.south) -- (argostg.north);
\draw[->,thick,green!40!black] (prod.south) -- (argoprod.north);

\end{tikzpicture}
%
}
\caption{Single-repo branch topology. The developer triggers \texttt{driveby-deploy.yml} to push manifests from \texttt{main} to a single \texttt{-next} branch; the Promoter creates PRs to active branches; quality gates control promotion; ArgoCD syncs from active branches. The BYOCI boundary is annotated: developer CI ends at \texttt{main}, XSDLC delivery begins at the deploy step.}
\label{fig:single-repo-branch-flow}
\end{figure}

The branch graph flows from \texttt{main} through the deploy workflow to a single \texttt{environment/<name>-next} branch per trigger. The Promoter creates promotion PRs from each \texttt{-next} branch to its active \texttt{environment/<name>} branch. Quality gates (when defined) evaluate before the PR is merged. ArgoCD watches each active branch and syncs the corresponding namespace.

% ============================================================================
\section{Multi-Environment Promotion}
\label{sec:multi-env-promotion}
% ============================================================================

The GitOps Promoter~\cite{gitops-promoter} automates environment promotion by watching for commit status signals and managing promotion pull requests. Figure~\ref{fig:multi-env-promotion} shows the three-environment model.

\begin{figure}[htbp]
\centering
\resizebox{\textwidth}{!}{%
% Multi-Environment Promotion — Chapter 6
% Shows: dev → staging → prod with quality gates between them
\begin{tikzpicture}[
    >={Stealth[length=2.5mm]},
    node distance=1.5cm,
    env/.style={
        rectangle,
        rounded corners=5pt,
        draw=#1!60!black,
        fill=#1!12,
        minimum width=3.5cm,
        minimum height=1.2cm,
        font=\small\bfseries,
        align=center
    },
    gate/.style={
        rectangle,
        rounded corners=3pt,
        draw=red!60!black,
        fill=red!10,
        minimum width=2.5cm,
        minimum height=0.7cm,
        font=\scriptsize,
        align=center
    },
    policy/.style={
        font=\tiny\itshape,
        text=#1!60!black,
        align=center
    }
]

% Environments
\node[env=blue] (dev) at (0,0) {Development\\(\texttt{dev})};
\node[env=orange] (staging) at (5.5,0) {Staging\\(\texttt{staging})};
\node[env=red] (prod) at (11,0) {Production\\(\texttt{prod})};

% Quality gates
\node[gate] (gate1) at (2.75,0) {DDT Quality\\Gate};
\node[gate] (gate2) at (8.25,0) {Manual\\Approval};

% Arrows
\draw[->,thick,blue!50!black] (dev) -- (gate1);
\draw[->,thick,orange!50!black] (gate1) -- (staging);
\draw[->,thick,orange!50!black] (staging) -- (gate2);
\draw[->,thick,red!50!black] (gate2) -- (prod);

% Policies
\node[policy=blue] at (0,-1.0) {autoMerge: true\\(auto-sync)};
\node[policy=orange] at (5.5,-1.0) {autoMerge: true\\(gated by DDT)};
\node[policy=red] at (11,-1.0) {autoMerge: false\\(manual approval)};

% Commit status flow
\draw[->,dashed,teal!50!black,thick] (gate1.north) -- ++(0,0.8) node[above,font=\tiny,teal!50!black] {CommitStatus: success/failure};

% Promoter label
\node[font=\scriptsize,text=black!60] at (5.5,1.5) {GitOps Promoter manages promotion PRs};

\end{tikzpicture}
%
}
\caption{Multi-environment promotion model. Development auto-syncs; staging promotion is gated by DDT validation (commit status); production requires manual approval.}
\label{fig:multi-env-promotion}
\end{figure}

% ----------------------------------------------------------------------------
\subsection{PromotionStrategy}
\label{subsec:promotion-strategy}
% ----------------------------------------------------------------------------

The PromotionStrategy custom resource defines the environment chain and the commit status keys that gate each transition. For the reference deployment, the strategy specifies three environments:

\begin{enumerate}
    \item \textbf{Development} (\texttt{dev} branch, \texttt{autoMerge: true}). Changes from the deploy workflow are automatically merged by the Promoter. No quality gate---development is the environment where new code lands first.
    \item \textbf{Staging} (\texttt{staging} branch, \texttt{autoMerge: true}). Promotion from development to staging is gated by the \texttt{staging-gate} commit status (the default key derived from the environment name). The Promoter creates a promotion PR, waits for the commit status to report success, and merges automatically. If the commit status reports failure, the PR remains open and promotion is blocked.
    \item \textbf{Production} (\texttt{prod} branch, \texttt{autoMerge: false}). Promotion from staging to production requires manual approval. A human reviewer must merge the promotion PR, providing a final checkpoint before production deployment.
\end{enumerate}

This two-tier model balances automation with safety. Development and staging promotions are fully automated (with DDT gating the staging transition); production promotion retains human oversight. Teams can adjust the \texttt{autoMerge} setting per environment to match their risk tolerance.

% ----------------------------------------------------------------------------
\subsection{CommitStatus as Quality Signal}
\label{subsec:commit-status-signal}
% ----------------------------------------------------------------------------

The integration between DriveBy and the GitOps Promoter is mediated by two commit status mechanisms:

\textbf{GitHub commit status.} The workflow's set-pending and report-success steps post commit status updates to the PR's head SHA via the GitHub API. The context key (e.g., \texttt{staging-gate}) is watched by the PromotionStrategy as the gate signal. This mechanism works for any GitHub-based workflow and does not require the GitOps Promoter.

\textbf{CommitStatus CRD.} When the GitOps Promoter is enabled, the workflow additionally creates a CommitStatus custom resource in the \texttt{driveby} namespace. The CommitStatus transitions through three phases: \texttt{pending} (set at workflow start), \texttt{success} (set after validation passes), or \texttt{failure} (set by the exit handler). The Promoter controller watches this resource and uses its phase to determine whether to merge the promotion PR. This CRD-based mechanism provides a Kubernetes-native alternative to the GitHub API and enables the Promoter to function even in environments where direct GitHub API access is restricted.

% ----------------------------------------------------------------------------
\subsection{ChangeTransferPolicy}
\label{subsec:change-transfer-policy}
% ----------------------------------------------------------------------------

The ChangeTransferPolicy defines how changes flow between environments at the Git level. For each environment, the policy maps a \textit{proposed} branch (e.g., \texttt{staging-next}) to an \textit{active} branch (e.g., \texttt{staging}). The Promoter creates pull requests from the proposed branch to the active branch and merges them when commit status gates pass. This branch-based model integrates naturally with Git-based review workflows: promotion PRs appear in the repository's PR list, carry commit status checks, and leave an auditable merge history.

% ----------------------------------------------------------------------------
\subsection{Promoter Reconciliation Loop}
\label{subsec:promoter-reconciliation}
% ----------------------------------------------------------------------------

The GitOps Promoter is itself a Kubernetes controller that implements the observe-diff-act pattern (Section~\ref{subsec:reconciliation-loop}) at the Git level rather than the cluster level.

\textbf{Observe.} The PromotionStrategy controller watches two resource types: GitRepository (which tracks the application repository) and CommitStatus (which reports quality gate results). It periodically polls the Git repository to detect changes on proposed branches and reads CommitStatus resources to determine gate outcomes.

\textbf{Diff.} For each environment in the promotion chain, the controller compares the proposed branch (\texttt{environment/<name>-next}) against the active branch (\texttt{environment/<name>}). If the branches differ, a promotion is pending. It then checks whether all required CommitStatuses report \texttt{success}. If any CommitStatus reports \texttt{failure} or \texttt{pending}, the promotion is blocked.

\textbf{Act.} When the proposed and active branches differ, the controller creates a promotion PR. When all commit status gates pass and \texttt{autoMerge} is \texttt{true}, the controller merges the PR automatically. When \texttt{autoMerge} is \texttt{false}, the PR remains open for human review.

The ChangeTransferPolicy is itself a second reconciliation layer: it is auto-created by the PromotionStrategy controller when the PromotionStrategy is applied. The PromotionStrategy reconciles the \textit{promotion logic} (which environments exist, what gates are required); the ChangeTransferPolicy reconciles the \textit{branch mechanics} (which branches map to which environments, when to create PRs). This two-layer design separates strategic configuration (the promotion chain) from tactical execution (the branch operations).

% ============================================================================
\section{Human-in-the-Loop Patterns}
\label{sec:human-in-the-loop}
% ============================================================================

The quality gate system is not fully autonomous. It automates objective assessments---DDT validation, health checks, load testing---while preserving human judgment for subjective assessments. This section maps the automation spectrum and identifies the decision points where human intervention is required or optional.

% ----------------------------------------------------------------------------
\subsection{The Automation Spectrum}
\label{subsec:automation-spectrum}
% ----------------------------------------------------------------------------

The system supports four points on the automation spectrum, controlled by two boolean parameters: \texttt{autoMerge} (on the PromotionStrategy, controlling whether the Promoter automatically merges promotion PRs) and \texttt{automated.selfHeal} (on the ArgoCD Application, controlling whether ArgoCD automatically syncs and self-heals).

\begin{enumerate}
    \item \textbf{Fully automated} (development): \texttt{autoMerge: true}, automated sync. After the developer triggers a deploy, changes flow from the \texttt{-next} branch, are auto-merged by the Promoter, and auto-synced by ArgoCD. No further human intervention after the initial trigger.
    \item \textbf{Automated with quality gate} (staging): \texttt{autoMerge: true}, automated sync, with a commit status gate. The Promoter waits for DDT validation to report success before merging. If validation fails, the PR remains open until the developer fixes the issue and re-pushes. Once the gate passes, promotion and sync are automatic.
    \item \textbf{Gated with human approval} (production): \texttt{autoMerge: false}, manual sync. The Promoter creates the PR and DDT validates the source environment, but a human must explicitly merge the promotion PR. Even after merge, ArgoCD requires manual sync (since automated sync is disabled). Two human checkpoints bookend the quality gate.
    \item \textbf{Fully manual}: no quality gate, manual merge, manual sync. This mode is possible but unused in the reference deployment---it offers no advantage over the gated model.
\end{enumerate}

% ----------------------------------------------------------------------------
\subsection{Decision Points}
\label{subsec:decision-points}
% ----------------------------------------------------------------------------

Figure~\ref{fig:human-in-the-loop} maps the decision points across the staging and production lanes.

\begin{figure}[htbp]
\centering
\resizebox{\textwidth}{!}{%
% Human-in-the-Loop — Chapter 6, Section 6.5.2
% Pipeline with green (automated) and amber (human) steps for staging vs prod
\begin{tikzpicture}[
    >={Stealth[length=2.5mm]},
    auto step/.style={
        rectangle,
        rounded corners=3pt,
        draw=green!60!black,
        fill=green!10,
        minimum width=2cm,
        minimum height=0.7cm,
        font=\scriptsize,
        align=center
    },
    human step/.style={
        rectangle,
        rounded corners=3pt,
        draw=orange!70!black,
        fill=orange!15,
        minimum width=2cm,
        minimum height=0.7cm,
        font=\scriptsize,
        align=center
    },
    lane label/.style={
        font=\small\bfseries,
        text=#1!70!black,
        anchor=east
    }
]

% Lane labels
\node[lane label=teal] at (-0.3,1.5) {Staging};
\node[lane label=violet] at (-0.3,-1.5) {Production};

% Staging lane (mostly automated)
\node[human step] (s1) at (1.5,1.5) {Code\\Author};
\node[auto step]  (s2) at (3.8,1.5) {Propagate};
\node[auto step]  (s3) at (6.1,1.5) {DDT Gate};
\node[auto step]  (s4) at (8.4,1.5) {Auto-merge};
\node[auto step]  (s5) at (10.7,1.5) {ArgoCD\\Sync};

\draw[->,thick] (s1) -- (s2);
\draw[->,thick] (s2) -- (s3);
\draw[->,thick] (s3) -- (s4);
\draw[->,thick] (s4) -- (s5);

% Production lane (human checkpoints)
\node[human step] (p1) at (1.5,-1.5) {Code\\Author};
\node[auto step]  (p2) at (3.8,-1.5) {Propagate};
\node[auto step]  (p3) at (6.1,-1.5) {Load Gate};
\node[human step] (p4) at (8.4,-1.5) {Manual\\Approve};
\node[human step] (p5) at (10.7,-1.5) {Manual\\Sync};

\draw[->,thick] (p1) -- (p2);
\draw[->,thick] (p2) -- (p3);
\draw[->,thick] (p3) -- (p4);
\draw[->,thick] (p4) -- (p5);

% Lane backgrounds
\begin{scope}[on background layer]
    \node[rectangle, rounded corners=4pt, fill=teal!5, draw=teal!20,
          fit={($(s1.north west)+(-0.2,0.3)$) ($(s5.south east)+(0.2,-0.3)$)}] {};
    \node[rectangle, rounded corners=4pt, fill=violet!5, draw=violet!20,
          fit={($(p1.north west)+(-0.2,0.3)$) ($(p5.south east)+(0.2,-0.3)$)}] {};
\end{scope}

% Legend
\node[auto step, minimum width=1.5cm] (leg1) at (3.5,-3.2) {Automated};
\node[human step, minimum width=1.5cm] (leg2) at (6.5,-3.2) {Human};
\node[font=\scriptsize] at (4.95,-3.2) {=};

% Annotation arrows for key decisions
\node[font=\tiny\itshape,text=green!50!black] at (8.4,0.6) {autoMerge: true};
\node[font=\tiny\itshape,text=orange!50!black] at (8.4,-0.6) {autoMerge: false};

\end{tikzpicture}
%
}
\caption{Human-in-the-loop patterns. Staging automates everything except code authorship; production adds human approval at the promotion PR merge and ArgoCD sync points. Green steps are automated; amber steps require human judgment.}
\label{fig:human-in-the-loop}
\end{figure}

All human intervention points in the system are:

\begin{enumerate}
    \item \textbf{Code authorship}: the developer writes the code and updates manifests. This is inherently human (or human-supervised, in the case of AI-assisted development).
    \item \textbf{PR review} (\texttt{autoMerge: false}): a human reviews the promotion PR and decides whether to merge. This checkpoint applies to production promotions, where the decision involves business timing, release coordination, and risk assessment---subjective factors that DDT does not measure.
    \item \textbf{ArgoCD sync} (manual sync): a human triggers synchronization in the ArgoCD UI. This checkpoint ensures that production deployments are initiated deliberately, not as an automatic consequence of a merge.
    \item \textbf{Failure diagnosis}: when a quality gate fails, a human (or autonomous agent) must diagnose the failure, fix the issue, and re-push. The DDT validation report provides per-principle diagnostics to guide this process.
\end{enumerate}

% ----------------------------------------------------------------------------
\subsection{Automation with Human Override}
\label{subsec:automation-override}
% ----------------------------------------------------------------------------

The design philosophy is: \textit{automate objective assessments, preserve human judgment for subjective assessments}. DDT validation, health checks, and load testing produce binary signals (pass/fail) based on measurable criteria. These are objective and automatable. Production readiness, business timing, and release coordination involve context that no automated system can capture. These are subjective and require human judgment.

The \texttt{autoMerge: false} setting is a deliberate architectural choice, not a limitation. It reflects the principle that the final decision to deploy to production should be an affirmative human act, even when all automated checks have passed. The quality gate provides confidence; the human provides authorization.

% ============================================================================
\section{Quality Gate as Control Loop}
\label{sec:control-loop}
% ============================================================================

This section synthesizes the architectural components described in Chapters~\ref{ch:architecture},~\ref{ch:kubernetes}, and the preceding sections of this chapter into a single control-theoretic framework. The quality gate system implements a \textit{closed-loop negative feedback} controller~\cite{astrom-feedback} that maintains API quality across the delivery pipeline.

% ----------------------------------------------------------------------------
\subsection{Control Theory Mapping}
\label{subsec:control-mapping}
% ----------------------------------------------------------------------------

The mapping from control theory to the DriveBy quality gate system is direct:

\begin{description}
    \item[Desired State (Setpoint):] The OpenAPI specification declares the desired API contract---the schemas, operations, error responses, security definitions, and documentation that the API should provide. This is the reference input against which the system measures deviation.

    \item[Plant:] The API deployment in the source environment (e.g., \texttt{perfect-api-dev}) is the controlled process. Its behavior---response schemas, error codes, latency characteristics---is the output being regulated.

    \item[Sensor:] The DriveBy CLI is the measurement instrument. It observes the plant's output by fetching the specification, validating it against the DDT principles, and producing a structured report. The sensor converts the plant's complex, multidimensional behavior into a scalar quality signal (pass/fail per principle, aggregate score).

    \item[Controller:] The Argo Workflow DAG is the decision-making component. It orchestrates the sensor (DriveBy validation steps), evaluates the results, and produces a control signal. The controller implements a simple threshold policy: if all checks pass, emit ``success''; if any check fails, emit ``failure''. The exit handler guarantees that the controller always produces a signal---no workflow terminates without an explicit quality verdict.

    \item[Actuator:] The CommitStatus CRD and the GitOps Promoter together form the actuator---the component that translates the control signal into physical action. A ``success'' signal causes the Promoter to merge the promotion PR, advancing the deployment. A ``failure'' signal causes the Promoter to hold the PR open, blocking advancement.

    \item[Feedback Loop:] The Promoter closes the loop. When promotion is blocked (failure signal), the developer must diagnose the issue, fix the code or specification, and re-trigger the deploy workflow. The new manifests update the \texttt{-next} branch, the Promoter updates the PR, and the quality gate re-evaluates. The system iterates until the plant's behavior converges with the desired state.
\end{description}

Figure~\ref{fig:quality-control-loop} illustrates this mapping.

\begin{figure}[htbp]
\centering
\resizebox{\textwidth}{!}{%
% Quality Gate as Control Loop — Chapter 6, Section 6.6
% Control theory block diagram: Setpoint (promotion policy) → Controller → Actuator → Plant (delivery system) → Sensor → feedback
% Mapping kept in lockstep with Table 6.1 (tab:control-mapping):
%   Setpoint = promotion-policy predicate (OpenAPI spec + XSDLC manifest)
%   Plant    = delivery system (pipeline + promotion workflow)
%   Actuator = CommitStatus + Promoter (merges/blocks promotion PR — acts ON the delivery system)
%   Sensor   = DriveBy CLI; measurement = validation verdict
\begin{tikzpicture}[
    >={Stealth[length=4mm,width=3mm]},
    block/.style={
        rectangle,
        rounded corners=4pt,
        draw=#1!60!black,
        fill=#1!12,
        minimum width=3cm,
        minimum height=1cm,
        font=\small,
        align=center
    },
    signal/.style={
        font=\footnotesize,
        text=gray!40!black
    },
    sum/.style={
        circle,
        draw=gray!60,
        fill=gray!5,
        minimum size=0.5cm,
        font=\scriptsize
    },
    flow/.style={->,line width=1.1pt}
]

% Setpoint: promotion-policy predicate
\node[block=violet] (setpoint) at (0,0) {\textbf{Setpoint}\\Promotion Policy\\(OpenAPI $+$ XSDLC)};

% Sum point
\node[sum] (sum) at (2.4,0) {$\Sigma$};

% Controller
\node[block=teal] (controller) at (4.7,0) {\textbf{Controller}\\Argo Workflow\\DAG};

% Actuator
\node[block=orange] (actuator) at (7.7,0) {\textbf{Actuator}\\CommitStatus +\\Promoter};

% Plant: the delivery system
\node[block=blue] (plant) at (10.7,0) {\textbf{Plant}\\Delivery System\\(pipeline $+$ promotion)};

% Sensor
\node[block=red!70!black] (sensor) at (7.7,-2.5) {\textbf{Sensor}\\DriveBy CLI\\(validation engine)};

% Forward path
\draw[flow,violet!50!black] (setpoint) -- (sum);
\draw[flow,gray!50!black] (sum) -- node[above,signal] {error} (controller);
\draw[flow,teal!50!black] (controller) -- node[above,signal] {pass/fail} (actuator);
\draw[flow,orange!50!black] (actuator) -- node[above,signal] {merge/block PR} (plant);

% Feedback path: the plant's deployed API (source env) is measured by the sensor
\draw[flow,blue!50!black] (plant.south) -- ++(0,-1) -- node[above,signal] {deployed API (source env)} (sensor.east);
\draw[flow,red!50!black] (sensor.west) -| node[left,signal,pos=0.3,align=right] {validation\\verdict} (sum.south);

% Negative feedback annotation
\node[font=\small\bfseries,text=red!60!black] at (2.8,-1.2) {$-$};

% Labels for control theory mapping
\node[font=\footnotesize\itshape,text=violet!50!black,anchor=north] at (0,-0.7) {(allowed transitions)};
\node[font=\footnotesize\itshape,text=teal!50!black,anchor=north] at (4.7,-0.7) {(orchestrator)};
\node[font=\footnotesize\itshape,text=orange!50!black,anchor=north] at (7.7,-0.7) {(acts on delivery system)};
\node[font=\footnotesize\itshape,text=blue!50!black,anchor=north] at (10.7,-0.7) {(promotion PRs, branches)};
\node[font=\footnotesize\itshape,text=red!50!black,anchor=north] at (7.7,-3.2) {(SDT validator)};

% Convergence annotation
\node[font=\footnotesize\itshape,text=gray!40!black,align=center] at (4.7,-3.9) {Closed-loop negative feedback: a non-conforming artifact blocks the transition\\$\rightarrow$ developer fixes $\rightarrow$ re-validates $\rightarrow$ only conforming artifacts are promoted};

\end{tikzpicture}
%
}
\caption{Quality gate as control loop. The OpenAPI specification is the setpoint; the DriveBy CLI is the sensor; the Argo Workflow is the controller; the CommitStatus and Promoter form the actuator. Negative feedback blocks promotion on validation failure, driving convergence.}
\label{fig:quality-control-loop}
\end{figure}

% ----------------------------------------------------------------------------
\subsection{Negative Feedback and Convergence}
\label{subsec:negative-feedback}
% ----------------------------------------------------------------------------

The feedback is \textit{negative}: validation failure produces a signal that opposes the direction of change (blocking promotion rather than advancing it). This is the defining characteristic of a regulating control loop~\cite{astrom-feedback}. Without negative feedback, a failing API would be promoted to staging and then to production, with quality problems detected only after deployment---the shift-right antipattern that DDT explicitly prevents.

The convergence dynamics depend on the developer response. Each quality gate evaluation is a discrete sampling of the plant's state. The developer's fix attempt modifies the plant. The next evaluation samples the modified state. If the fix addresses the failing principle, the quality signal transitions from ``failure'' to ``success'' and the loop converges. If the fix is incomplete, the signal remains ``failure'' and the loop iterates.

The level-triggered nature of the Kubernetes controller pattern (Section~\ref{subsec:reconciliation-loop}) ensures that the control loop is robust to transient failures. If a workflow fails due to infrastructure issues (pod eviction, network timeout) rather than validation failure, the exit handler posts a ``failure'' status. The developer re-pushes (or the Promoter detects the branch change), and the workflow re-executes. The system does not require manual re-triggering because the Promoter continuously reconciles the branch state.

% ----------------------------------------------------------------------------
\subsection{Connecting Theory to Implementation}
\label{subsec:theory-implementation}
% ----------------------------------------------------------------------------

This control-theoretic framing connects the abstract ontological state reconciliation pattern introduced in Section~\ref{subsec:ontological-reconciliation} of Chapter~\ref{ch:related-work} to the concrete implementation described in this thesis.

In Section~\ref{subsec:ontological-reconciliation}, we defined ontological state reconciliation as requiring three components: an ontological artifact (the specification), a reconciliation agent (the validator), and a correction mechanism (the quality gate). The control loop formalization adds precision: the specification is the setpoint, the validator is the sensor, and the quality gate is the controller-actuator pair that closes the feedback loop.

The three DDT axioms map directly to control-theoretic requirements. The \textit{Axiom of Completeness} requires that the setpoint (specification) is expressive enough to fully characterize the desired behavior---an incomplete setpoint produces an under-determined control loop that cannot distinguish conforming from non-conforming behavior. The \textit{Axiom of Determinism} requires that the sensor (validator) produces consistent measurements---a non-deterministic sensor would cause the control loop to oscillate between pass and fail for identical plant states. The \textit{Axiom of Observability} requires that the sensor's output is structured and actionable---an unobservable control loop cannot inform corrective action, and the feedback degenerates into a binary gate with no diagnostic value.

The quality gate system thus instantiates a well-defined control loop in which every component---setpoint, sensor, controller, actuator, and feedback path---has a concrete implementation in the DriveBy architecture. This is not a metaphorical mapping; it is a structural correspondence that explains why the system converges, why it is robust to transient failures, and why the axiomatic foundation is necessary rather than merely desirable.

% ============================================================================
\section{CI/CD of DriveBy Itself}
\label{sec:cicd-driveby}
% ============================================================================

DriveBy's own development is governed by three GitHub Actions~\cite{github-actions} workflows that implement continuous integration, continuous delivery, and release automation. These pipelines demonstrate that the project practices what it preaches: automated quality gates at every stage of the delivery lifecycle.

% ----------------------------------------------------------------------------
\subsection{Continuous Integration}
\label{subsec:ci-pipeline}
% ----------------------------------------------------------------------------

The \texttt{ci.yml} workflow triggers on every push to \texttt{main} and every pull request. It runs a single \texttt{build-and-test} job on Ubuntu with Go~1.21 that executes four steps:

\begin{enumerate}
    \item \textbf{Build}: Compiles the DriveBy binary with version and commit SHA injected via \texttt{ldflags}.
    \item \textbf{Static analysis}: Runs \texttt{go vet} to detect common Go errors.
    \item \textbf{Test}: Executes the test suite from \texttt{driveby-cli/test/} with coverage profiling (\texttt{-coverprofile=coverage.out}).
    \item \textbf{Documentation check}: Runs \texttt{tools/check-docs.sh} to verify that all documented components exist and all source files have corresponding documentation entries.
\end{enumerate}

The documentation check is a meta-validation step: it ensures that the CLAUDE.md files and documentation artifacts remain synchronized with the codebase. This check embodies a DDT-adjacent principle---documentation completeness should be verified automatically, not assumed.

On pull requests, the workflow uploads the coverage artifact for review. On pushes to \texttt{main}, a subsequent \texttt{docker} job builds the container image (without pushing) to verify the Docker build remains functional.

% ----------------------------------------------------------------------------
\subsection{Docker Image Publishing}
\label{subsec:docker-publish}
% ----------------------------------------------------------------------------

The \texttt{docker-publish.yml} workflow publishes container images to the GitHub Container Registry (GHCR) on pushes to \texttt{main} when files in \texttt{driveby-cli/} change. It uses Docker Buildx~\cite{docker-buildx} for multi-platform builds, injecting three build arguments:

\begin{itemize}
    \item \texttt{VERSION}: The current Git tag or \texttt{dev} for unreleased commits.
    \item \texttt{COMMIT}: The short Git SHA.
    \item \texttt{BUILD\_DATE}: The ISO-8601 build timestamp.
\end{itemize}

Images are tagged with \texttt{latest} and the short commit SHA. Layer caching via GitHub Actions Cache reduces build times for incremental changes. The resulting image at \texttt{ghcr.io/meter-peter/driveby} is the same image used by the Argo Workflows validation pipeline.

% ----------------------------------------------------------------------------
\subsection{Release Pipeline}
\label{subsec:release-pipeline}
% ----------------------------------------------------------------------------

The \texttt{release.yml} workflow triggers on semantic version tag pushes (\texttt{v*}) and orchestrates the full release process. Figure~\ref{fig:release-pipeline} shows the pipeline structure.

\begin{figure}[htbp]
\centering
\resizebox{\textwidth}{!}{%
% Release Pipeline — Chapter 6
% Shows: Tag push → test gate → goreleaser + docker + helm (parallel fan-out)
\begin{tikzpicture}[
    >={Stealth[length=2.5mm]},
    node distance=0.9cm,
    stage/.style={
        rectangle,
        rounded corners=3pt,
        draw=#1!60!black,
        fill=#1!12,
        minimum width=3cm,
        minimum height=0.8cm,
        font=\small,
        align=center
    },
    gate/.style={
        diamond,
        draw=red!60!black,
        fill=red!10,
        minimum width=1.5cm,
        minimum height=1cm,
        font=\small,
        align=center,
        inner sep=1pt
    },
    artifact/.style={
        rectangle,
        rounded corners=3pt,
        draw=#1!50!black,
        fill=#1!8,
        minimum width=2.8cm,
        minimum height=0.6cm,
        font=\scriptsize,
        align=center
    }
]

% Trigger
\node[stage=violet] (tag) at (0,0) {Tag Push\\(\texttt{v*})};

% Test gate
\node[gate] (test) at (3,0) {test};
\draw[->,thick] (tag) -- (test);

% Fan-out (three parallel jobs)
\node[stage=blue] (goreleaser) at (7,1.5) {GoReleaser};
\node[stage=teal] (docker) at (7,0) {Docker Build};
\node[stage=orange] (helm) at (7,-1.5) {Helm Package};

\draw[->,thick] (test) -- ++(1.5,0) |- (goreleaser);
\draw[->,thick] (test) -- (docker);
\draw[->,thick] (test) -- ++(1.5,0) |- (helm);

% Fan-out label
\node[font=\tiny\itshape,text=black!50] at (5,0.8) {parallel};

% Artifacts
\node[artifact=blue] (bins) at (11,2.0) {6 platform binaries\\+ SHA256 checksums};
\node[artifact=blue] (release) at (11,1.0) {GitHub Release};
\node[artifact=teal] (image) at (11,0) {ghcr.io/.../driveby\\:3.0.0, :3.0, :latest};
\node[artifact=orange] (chart) at (11,-1.5) {oci://ghcr.io/.../\\charts/driveby:3.0.0};

\draw[->,thick,blue!50!black] (goreleaser) -- (bins);
\draw[->,thick,blue!50!black] (goreleaser) -- (release);
\draw[->,thick,teal!50!black] (docker) -- (image);
\draw[->,thick,orange!50!black] (helm) -- (chart);

\end{tikzpicture}
%
}
\caption{Release pipeline triggered by a version tag push. After the test gate passes, three parallel jobs produce platform binaries, Docker images, and a Helm chart.}
\label{fig:release-pipeline}
\end{figure}

The pipeline runs four jobs:

\textbf{Test gate.} Re-runs \texttt{go vet} and the test suite. All subsequent jobs depend on this gate---no artifacts are produced if tests fail.

\textbf{GoReleaser.} Uses GoReleaser~\cite{goreleaser}~v2 to produce cross-compiled binaries for six platforms: Linux, macOS, and Windows on both amd64 and arm64 architectures. GoReleaser generates SHA256 checksums for each binary and creates a GitHub Release with the binaries, checksums, and an auto-generated changelog attached. The release changelog summarizes commits since the previous tag.

\textbf{Docker.} Builds and pushes the container image with three semver-derived tags: the full version (\texttt{0.4.0}), the major.minor (\texttt{0.4}), and \texttt{latest}. This tagging strategy enables consumers to pin to a specific patch version, track a minor version, or always use the latest release.

\textbf{Helm.} Packages the Helm chart from \texttt{kubernetes/helm/driveby/} and pushes it to the OCI registry at \texttt{oci://ghcr.io/meter-peter/charts/driveby}. The chart version matches the release tag, ensuring that the chart, container image, and CLI binary are version-aligned.

The three artifact-producing jobs run in parallel after the test gate, minimizing release time. A typical release completes in under five minutes from tag push to all artifacts published.

% ============================================================================
\section{Operational Deployment}
\label{sec:operational-deployment}
% ============================================================================

The quality gate system is deployed on \texttt{private.novelcore.org}, a Kubernetes cluster managed by Kubecore with ArgoCD~\cite{argocd}, Crossplane~\cite{crossplane}, Argo Events~\cite{argo-events}, Argo Workflows~\cite{argo-workflows}, Traefik, and cert-manager pre-installed.

% ----------------------------------------------------------------------------
\subsection{Prerequisites}
\label{subsec:prerequisites}
% ----------------------------------------------------------------------------

The deployment assumes six cluster components:

\begin{enumerate}
    \item \textbf{Crossplane} (v1.14+) with the \texttt{provider-kubernetes} provider for managing in-cluster resources.
    \item \textbf{ArgoCD} (v2.8+) for GitOps-based application deployment across dev, staging, and prod namespaces.
    \item \textbf{Argo Workflows} (v3.5+) for container-native workflow execution.
    \item \textbf{Argo Events} (v1.9+) for event-driven workflow triggering.
    \item \textbf{Traefik} (v2.10+) as the Ingress controller for routing webhook traffic.
    \item \textbf{cert-manager} (v1.13+) for automated TLS certificate provisioning.
\end{enumerate}

These components are managed independently of the DriveBy installation and are assumed to be healthy. The DriveBy Helm chart creates resources that reference these components (Ingress resources reference the Traefik ingress class, certificate annotations reference the cert-manager issuer) but does not install them.

% ----------------------------------------------------------------------------
\subsection{Verification}
\label{subsec:verification}
% ----------------------------------------------------------------------------

After deployment, the operator verifies the system through four checks:

\begin{enumerate}
    \item \textbf{Crossplane readiness}: \texttt{kubectl get composite} confirms that XSDLC resources report \texttt{READY=True}.
    \item \textbf{Argo Events health}: \texttt{kubectl get eventbus,eventsource,sensor -n driveby} confirms all three resources are in a ready state.
    \item \textbf{Webhook connectivity}: A test webhook payload sent to the Ingress endpoint returns a 200 response.
    \item \textbf{End-to-end validation}: Opening a test PR on the monitored repository triggers a workflow run. The operator watches \texttt{kubectl get workflows -n driveby} for workflow completion and verifies that commit status and PR comments appear on the test PR.
\end{enumerate}

% ============================================================================
\section{Summary}
\label{sec:pipeline-summary}
% ============================================================================

This chapter described the operational GitOps pipeline that connects DriveBy's validation capabilities to the software delivery lifecycle.

The promotion-driven quality gate (Section~\ref{sec:promotion-driven-gate}) validates the source environment before allowing promotion, implementing the shift-left principle at the deployment boundary. The event-driven pipeline (Section~\ref{sec:event-pipeline-detail}) transforms GitHub webhooks into structured validation workflows with sub-90-second latency. The single-repo architecture (Section~\ref{sec:single-repo}) eliminates config-repo duplication and establishes the BYOCI principle: XSDLC owns delivery (promotion and quality gates) while developers retain full control over their CI pipeline.

The multi-environment promotion model (Section~\ref{sec:multi-env-promotion}) automates the dev $\rightarrow$ staging transition through DDT-gated commit status signals while retaining manual approval for production. The Promoter itself is a reconciliation controller (Section~\ref{subsec:promoter-reconciliation}), implementing observe-diff-act at the Git level. Human-in-the-loop patterns (Section~\ref{sec:human-in-the-loop}) define a clear automation spectrum: objective assessments (validation, load testing) are automated; subjective assessments (production readiness, business timing) are preserved for human judgment.

The quality gate as control loop (Section~\ref{sec:control-loop}) synthesizes the entire system into a control-theoretic framework, mapping the OpenAPI specification to a setpoint, the DriveBy CLI to a sensor, the Argo Workflow to a controller, and the Promoter to an actuator. This closed-loop negative feedback system ensures convergence: validation failure blocks promotion, driving developers to fix issues until the API conforms to its specification. The three DDT axioms---Completeness, Determinism, and Observability---are shown to be necessary conditions for the control loop to function correctly, connecting the theoretical foundations of Chapter~\ref{ch:related-work} (Section~\ref{subsec:ontological-reconciliation}) to the concrete implementation.

The CI/CD pipelines for DriveBy itself (Section~\ref{sec:cicd-driveby}) demonstrate end-to-end release automation: from test gate through parallel artifact production (binaries, images, Helm chart) to registry publication. Together, these sections show that DDT validation integrates naturally into modern GitOps workflows---not as an afterthought bolted onto an existing pipeline, but as a first-class quality gate in a declarative, event-driven, control-theoretic delivery system.

Chapter~\ref{ch:evaluation} evaluates this system empirically, measuring DDT's effectiveness across controlled and large-scale settings.
